\section{Contraintes, hypothèses et décisions d'architecture}

\subsection{Contraintes}
\label{sec:contraintes}

Ces contraintes ne se négocient pas~: la conception s'y plie.

{\small
\begin{longtable}{@{}l>{\raggedright\arraybackslash}p{70mm}>{\raggedright\arraybackslash}p{72mm}@{}}
    \toprule
    \textbf{Réf.} & \textbf{Contrainte} & \textbf{Conséquence pour la conception} \\
    \midrule\endhead
    \bottomrule\endfoot
    C1 & Quota de 100 requêtes par jour pour un compte standard [E2]. Chaque réponse porte un en-tête \texttt{RateLimit} (restant, secondes avant renouvellement) et \texttt{RateLimit-Policy} (quota). Les écritures comptent [E4]. & Un composant unique compte, cache et limite toutes les requêtes (§\ref{sec:quota}). Le quota réel est lu dans l'en-tête, non supposé. \\
    C2 & L'API de Tado est non officielle pour les tiers~; la spécification communautaire [E3] signale que les chemins de zone ne valent que pour les installations \emph{classiques}, pas pour Tado~X. & Toute dépendance à un chemin est isolée dans un seul composant (\texttt{J-03}). Hypothèse H1. \\
    C3 & Authentification par code d'appareil [E1, N3]~: jeton d'accès de courte vie, jeton de rafraîchissement \emph{à rotation} (chaque usage invalide le précédent), valable 30 jours au plus. & Le jeton le plus récent est toujours écrit avant d'être utilisé~; une montre inutilisée plus de 30 jours impose une nouvelle liaison (\texttt{TPB-CPT-005}). \\
    C4 & La géolocalisation de Tado et la restitution du mode \emph{Auto} (\texttt{DELETE presenceLock}) exigent un abonnement Auto-Assist~[E3]. & Le mode \emph{Auto} n'est offert que s'il est disponible (H2). \\
    C5 & Le pkjs ne s'exécute que pendant que la watchapp est ouverte, ou qu'une page de configuration l'est. Il n'existe pas de tâche de fond. & Pas d'alerte poussée sans ouvrir l'application (DEC-08). \\
    C6 & Les tampons AppMessage sont de taille fixe~; un message plus gros que le tampon est perdu. & Messages petits, paginés (§\ref{sec:protocole}). \\
    C7 & Stockage de la watchapp~: 4~Ko sur aplite, basalt, chalk et diorite~; 1~Mo sur emery, gabbro et flint [E5]. Mémoire de travail réduite sur aplite. & La watchapp ne conserve que la dernière situation (quelques centaines d'octets) (\texttt{TPB-INF-001}). \\
    C8 & Pas de clavier, quatre boutons, écrans de 144\,×\,168 à 260\,×\,260. & Pas de saisie de texte (\texttt{TPB-UTI-003})~; un moteur de listes unique pour toutes les formes d'écran. \\
\end{longtable}}

\subsection{Hypothèses}
\label{sec:hypotheses}

\begin{description}
  \item[H1. Installation classique.] Confirmé par l'utilisateur~: la maison est équipée du
    système Tado classique, non de Tado~X. Si cela changeait, les chemins de zone changeraient
    d'hôte et de forme~: seul \texttt{J-03} serait à réécrire.
  \item[H2. Pas d'Auto-Assist.] L'utilisateur n'a pas l'abonnement Auto-Assist, mais la
  géolocalisation fonctionne dans son application Android Tado. Le quota est donc celui d'un compte
  standard (100 requêtes par jour, à confirmer par l'en-tête). Que l'API accepte de rendre la
  présence à la géolocalisation (\texttt{DELETE presenceLock}) sans l'abonnement n'est pas établi~:
  l'application ne le suppose pas, elle découvre le refus de Tado (§\ref{sec:presence}).
  \item[H3. Téléphone à portée.] La montre n'atteint Tado que par le téléphone associé
    (cahier des charges, \texttt{TPB-ACT-009}).
  \item[H4. Stockage local du pkjs.] Le \texttt{localStorage} du pkjs survit entre deux
    ouvertures de la watchapp. Il peut être vidé par l'utilisateur (désinstallation, nettoyage)~:
    l'application doit alors redemander la liaison.
  \item[H5. Fuseau et unités de la maison.] La date locale de la maison (pour demander le
    rapport du jour) et l'unité de température viennent des réglages de la maison
    (\texttt{GET /homes/\{id\}}), non de la montre.
\end{description}

\subsection{Vue d'ensemble}

\begin{figure}[H]
  \centering
  \begin{tikzpicture}[
      font=\sffamily\footnotesize,
      bloc/.style={draw=exigBleu, fill=exigBleu!5, rounded corners=3pt, align=center,
                   text width=38mm, minimum height=22mm},
      ext/.style={draw=exigGris, fill=exigGris!6, rounded corners=2pt, align=center,
                  text width=34mm, minimum height=14mm},
      fleche/.style={Stealth-Stealth, draw=exigGris!90!black},
      lbl/.style={font=\sffamily\scriptsize, text=exigBleu!80!black, align=center},
    ]
    \node[bloc] (mon) at (0, 0)    {\textbf{Watchapp}\\langage C\\écrans, boutons,\\dernière situation};
    \node[bloc] (tel) at (6.4, 0)  {\textbf{pkjs}\\JavaScript, téléphone\\logique Tado, cache,\\quota, vues};
    \node[ext]  (log) at (12.2, 1.2)  {login.tado.com\\\textit{autorisation}};
    \node[ext]  (api) at (12.2, -1.2) {my.tado.com\\\textit{API v2}};
    \draw[fleche] (mon) -- node[lbl, above=1pt]{AppMessage} node[lbl, below=1pt]{Bluetooth} (tel);
    \draw[fleche] (tel) -- node[lbl, above=1pt]{HTTPS} (log);
    \draw[fleche] (tel) -- node[lbl, below=1pt]{HTTPS} (api);
  \end{tikzpicture}
  \caption{Contexte~: deux programmes, deux services de Tado}
  \label{fig:contexte}
\end{figure}

La watchapp ne parle jamais à Tado. Elle envoie des \emph{demandes} au pkjs et reçoit des
\emph{vues}~; le pkjs porte toute la connaissance de Tado, du quota et des jetons. Aucun
serveur n'est ajouté entre le téléphone et Tado.

\subsection{Décisions d'architecture}
\label{sec:decisions}

\begin{decision}{DEC-01}{Watchapp en C et pkjs en JavaScript}
  \attr{Contexte} Une montre Pebble n'a pas de réseau propre~; le SDK offre le C natif (toutes
  plateformes), le pkjs pour le téléphone, et des environnements JavaScript plus récents sur la montre.
  \attr{Décision} La watchapp est écrite en C et le pkjs en JavaScript, comme les applications
  précédentes de l'utilisateur (cohérence des outils et du savoir-faire).
  \attr{Alternatives écartées} JavaScript sur la montre (Alloy)~: moins éprouvé, et le support
  des sept plateformes est à établir. Rocky.js~: réservé aux cadrans.
  \attr{Conséquences} Deux langages, un protocole à tenir à jour (§\ref{sec:protocole}).
  \attr{Exigences} \refexig{TPB-PLT-001}, \refexig{TPB-PLT-002}.
\end{decision}

\begin{decision}{DEC-02}{Toute la logique Tado sur le téléphone}
  \attr{Contexte} Il faut limiter la mémoire de la montre et ne pas dupliquer des règles.
  \attr{Décision} Le pkjs fait toutes les requêtes, tous les calculs (intensité, niveaux,
  abréviations, choix des raccourcis à proposer) et ne transmet à la montre que des vues
  prêtes à afficher. La watchapp n'interprète aucune donnée de Tado.
  \attr{Alternatives écartées} Faire calculer la montre~: code dupliqué en C, mémoire sur aplite.
  \attr{Conséquences} Les règles se testent en JavaScript sans montre. La montre est
  inutile sans téléphone, ce qui est déjà vrai (H3).
  \attr{Exigences} \refexig{TPB-PIE-002}, \refexig{TPB-HIS-002}, \refexig{TPB-RAC-003} à
  \refexig{TPB-RAC-005}, \refexig{TPB-PLT-006}.
\end{decision}

\begin{decision}{DEC-03}{Liaison par code d'appareil, jetons sur le téléphone}
  \attr{Contexte} Tado a retiré l'authentification par mot de passe au profit du code
  d'appareil [E1]~; la montre n'a pas de clavier.
  \attr{Décision} Le flux de la RFC~8628 est conduit par le pkjs. La montre affiche le code
  à saisir sur l'adresse de validation, que le téléphone ouvre (page de configuration ou
  navigateur). Le pkjs interroge Tado toutes les cinq secondes jusqu'à validation ou
  expiration (300~s). Les jetons sont conservés dans le \texttt{localStorage} du pkjs~; le
  mot de passe ne passe jamais par l'application.
  \attr{Alternatives écartées} Saisie du mot de passe dans une page de configuration~: flux
  retiré par Tado, mot de passe conservé.
  \attr{Conséquences} Un écran de liaison à ajouter aux maquettes (§\ref{sec:maquettes}).
  Le jeton de rafraîchissement est réécrit à chaque usage (C3).
  \attr{Exigences} \refexig{TPB-CPT-001}, \refexig{TPB-CPT-003} à \refexig{TPB-CPT-005},
  \refexig{TPB-SEC-001}, \refexig{TPB-SEC-002}.
\end{decision}

\begin{decision}{DEC-04}{Requêtes groupées}
  \attr{Contexte} Chaque requête coûte du quota (C1).
  \attr{Décision} L'état de toutes les pièces vient d'une seule requête (\texttt{zoneStates}), la
  présence d'une seule (\texttt{state}). Les raccourcis utilisent les routes groupées
  \texttt{POST /overlay} et \texttt{DELETE /overlay?rooms=}. Aucune requête par pièce n'est faite pour la liste.
  \attr{Conséquences} Le coût d'ouverture est de deux requêtes, quel que soit le nombre de pièces.
  \attr{Exigences} \refexig{TPB-ACT-005}, \refexig{TPB-RAC-002}, \refexig{TPB-RAC-003}.
\end{decision}

\begin{decision}{DEC-05}{Historique par rapport du jour, avec cache}
  \attr{Contexte} L'historique doit être complet même montre fermée (\refexig{TPB-HIS-007}).
  \attr{Décision} Le pkjs demande le rapport du jour de chaque pièce (\texttt{dayReport}), qui
  contient les périodes d'appel de chaleur, et en tire douze colonnes. Un rapport de jour échu
  est conservé définitivement~; celui du jour courant se renouvelle au plus toutes les 30 minutes.
  \attr{Alternatives écartées} Échantillonner la puissance à intervalles réguliers~: consomme du
  quota à chaque ouverture, laisse des trous lorsque l'application est fermée.
  \attr{Conséquences} Un historique complet coûte une requête par pièce, deux par pièce avant
  midi (la fenêtre de douze heures chevauche minuit). Les niveaux de chauffe du rapport sont au
  nombre de quatre, non des pourcentages~: voir DEC-11 et le retour R3 (§\ref{sec:retours}).
  \attr{Exigences} \refexig{TPB-HIS-001}, \refexig{TPB-HIS-002}, \refexig{TPB-HIS-007}, \refexig{TPB-HIS-008}.
\end{decision}

\begin{decision}{DEC-06}{Un seul gardien du quota}
  \attr{Contexte} Les consultations ne doivent pas empêcher une commande (\refexig{TPB-ACT-006}).
  \attr{Décision} Toutes les requêtes passent par le composant \texttt{J-05}. Il lit les en-têtes
  \texttt{RateLimit}, refuse toute \emph{lecture} lorsqu'il reste dix requêtes ou moins, et
  n'accepte plus que les commandes. Il sert les données du cache lorsqu'il refuse.
  \attr{Conséquences} Lorsque le quota est épuisé, l'application l'annonce
  (\refexig{TPB-ACT-007}) et affiche le cache avec son âge. Le nombre de requêtes restantes n'est
  jamais montré (\refexig{TPB-ACT-010}).
  \attr{Exigences} \refexig{TPB-ACT-005} à \refexig{TPB-ACT-007}, \refexig{TPB-ACT-010}.
\end{decision}

\begin{decision}{DEC-07}{Dernière situation à deux endroits}
  \attr{Contexte} L'ouverture doit être immédiate (\refexig{TPB-ACT-001}), y compris sans téléphone.
  \attr{Décision} La watchapp conserve dans son stockage la dernière vue de la liste, avec
  l'heure de son obtention~; le pkjs conserve dans son cache toutes les données de Tado
  avec leur âge. À l'ouverture, la watchapp affiche sa copie, puis la remplace.
  \attr{Conséquences} Quelques centaines d'octets sur la montre, compatibles avec 4~Ko.
  \attr{Exigences} \refexig{TPB-ACT-001}, \refexig{TPB-ACT-002}, \refexig{TPB-ACT-009},
  \refexig{TPB-INF-001}, \refexig{TPB-INF-002}.
\end{decision}

\begin{decision}{DEC-08}{Pas de serveur, pas de tâche de fond}
  \attr{Contexte} Le pkjs ne tourne pas lorsque la watchapp est fermée (C5).
  \attr{Décision} L'application n'a ni serveur propre ni exécution en arrière-plan. Elle ne peut donc
  pas prévenir l'utilisateur d'une alerte qu'il n'a pas cherchée.
  \attr{Alternatives écartées} Un serveur qui interroge Tado et pousse des notifications~:
  hébergement, jetons hors du téléphone, et quota consommé en continu.
  \attr{Conséquences} \refexig{TPB-ALE-006} (Souhaitable) n'est pas satisfaite dans la première
  version~; voir le retour R4.
  \attr{Exigences} \refexig{TPB-ALE-006}, \refexig{TPB-INF-003}.
\end{decision}

\begin{decision}{DEC-09}{Un seul code C pour sept plateformes}
  \attr{Contexte} Les plateformes diffèrent par la forme, la couleur et la taille.
  \attr{Décision} Un seul code, adapté à la compilation par les macros du SDK (\texttt{PBL\_COLOR},
  \texttt{PBL\_ROUND}, \texttt{PBL\_DISPLAY\_WIDTH}) et par une table de constantes par plateforme
  (hauteurs de ligne, nombre de colonnes d'historique, largeur de la barre d'actions).
  \attr{Conséquences} Les maquettes de chaque plateforme (§\ref{sec:maquettes}) fournissent les
  valeurs de la table.
  \attr{Exigences} \refexig{TPB-PLT-001} à \refexig{TPB-PLT-007}, \refexig{TPB-AFF-003}.
\end{decision}

\begin{decision}{DEC-10}{Une commande n'est jamais présumée réussie}
  \attr{Contexte} Afficher la valeur demandée avant la confirmation de Tado tromperait l'utilisateur.
  \attr{Décision} Après chaque commande, le pkjs relit l'état de la pièce (ou de la maison) et
  ne renvoie à la montre que cet état. Une commande refusée est annoncée comme telle, sans
  réessai automatique, et sans file d'attente (\refexig{TPB-PIL-011}).
  \attr{Conséquences} Une commande coûte deux requêtes (l'écriture et la relecture).
  \attr{Exigences} \refexig{TPB-PIL-008} à \refexig{TPB-PIL-011}, \refexig{TPB-ACT-008}.
\end{decision}

\begin{decision}{DEC-11}{Intensité de chauffe pondérée}
  \attr{Contexte} Le cahier des charges définit l'intensité d'une heure comme la puissance moyenne
  de l'heure, en pourcentage de la pleine puissance (\refexig{TPB-HIS-002}). Le rapport du jour ne
  donne pas de pourcentage mais quatre niveaux (aucun, faible, moyen, fort) par période.
  \attr{Décision} Chaque niveau reçoit un poids~: aucun 0~\%, faible 25~\%, moyen 50~\%, fort
  100~\%. L'intensité d'une heure est la moyenne de ces poids, pondérée par la durée de chaque période
  dans l'heure. Elle est rangée dans les quatre plages du cahier des charges.
  \attr{Conséquences} Les poids sont provisoires~: à étalonner en comparant, sur quelques jours,
  les niveaux du rapport à la puissance en pourcentage de \texttt{zoneStates}.
  \attr{Exigences} \refexig{TPB-HIS-001}, \refexig{TPB-HIS-002}.
\end{decision}

\begin{decision}{DEC-12}{Textes fixes sur la montre, textes variables sur le téléphone}
  \attr{Contexte} Deux langues (français et anglais) et deux unités de température
  (\refexig{TPB-AFF-005}, \refexig{TPB-AFF-006}).
  \attr{Décision} Les libellés fixes (titres, modes, boutons) sont des ressources de la
  watchapp, une par langue, choisies selon la langue de la montre. Le pkjs envoie des codes
  pour ces libellés et des textes pour ce qui varie (noms de pièces, abréviations). Les
  températures voyagent en dixièmes de degré dans l'unité du compte~; la watchapp
  les met en forme (virgule ou point).
  \attr{Conséquences} Ajouter une langue ne touche que les ressources de la montre.
  \attr{Exigences} \refexig{TPB-AFF-005}, \refexig{TPB-AFF-006}.
\end{decision}
